\documentclass[runningheads]{llncs}
%
\usepackage[T1]{fontenc}
\usepackage{graphicx}
\usepackage{amsmath}
\usepackage{amssymb}
\usepackage{booktabs}
\usepackage{url}
\usepackage{algorithm}
\usepackage{algorithmic}

\begin{document}
%
\title{Digital Processing Pipeline for Integrity Restoration in Pixel Art Generated by Diffusion Models}
%
\titlerunning{Digital Processing Pipeline for Pixel Art Integrity Restoration}
%
\author{Nilton Santos\inst{1} \and
João Pena\inst{1} \and
Renan Vinicius Aranha\inst{1}}
%
\authorrunning{N. Santos, J. Pena, and R. V. Aranha}
%
\institute{Universidade Federal de Mato Grosso (UFMT), Cuiabá, Brazil\\
\email{\{nilton.santos, joao.pena, renan.aranha\}@ufmt.br}}
%
\maketitle              % typeset the header of the contribution
%
\begin{abstract}
Digital asset production in game development depends on generative models that rarely satisfy technical constraints exactly. Even when diffusion-based architectures are mature, random noise, interpolation artifacts, and mixels introduce inconsistencies that compromise grid alignment, palette integrity, and engine compatibility. This article presents a digital processing pipeline for Pixel Art integrity restoration, providing a deterministic bridge between synthetic image generation and production-ready sprites. The contribution of the work is not only algorithmic, but architectural: the system combines a nearest-neighbor resampling stage for grid reconstruction, a K-Means quantization layer for chromatic reduction, and an alpha binarization module for silhouette normalization. The article reorganizes the broader undergraduate thesis into a systems-oriented paper, emphasizing mathematical formulation, design rationale, implementation decisions, and prototype capabilities. The resulting platform demonstrates how classical image processing methods can be operationalized in an automated environment that makes generative workflows more inspectable, reproducible, and available to engineering contexts that require accessible digital instrumentation for game assets.

\keywords{Generative AI \and Computer Vision \and Game Assets \and Pixel Art \and Digital Processing Pipeline.}
\end{abstract}
%
%
%
\section{Introduction}

Pixel Art is a discipline governed by severe mathematical constraints: each pixel must inhabit a discrete coordinate on a grid and belong to a limited chromatic palette \cite{mercioni2024innovative,zufri2022}. When transposing diffusion-based generative AI models to this domain, a fundamental technical conflict emerges: these models operate by noise reduction in continuous spaces, producing interpolations and anti-aliasing that are often opposed to the indexed purity required by the style \cite{chen2023pixart,ribeiro2024image}. Recent surveys on generative AI-driven game asset production highlight that while visual quality has improved, technical compliance remains a bottleneck for direct pipeline integration \cite{kim2025survey}. In practice, many outputs of these systems result in technically problematic assets, presenting \textit{mixels} --- elements that float between grid cells --- and chromatic noise that prevents direct use in game engines \cite{lowenstrom2024guiding,coutinho2022}. This problem is not merely aesthetic. In production settings, especially in independent game development, assets must satisfy strict technical requirements regarding palette control, transparency, scaling behavior, and predictability under animation or engine import. An image that only ``resembles'' Pixel Art at first glance may still fail as an operational sprite (Fig.~\ref{fig:artifacts}).

\begin{figure}[ht]
\centering
\includegraphics[width=0.45\linewidth]{figures/ai_pixel_artifacts.png}
\caption{Recurring artifacts in raw diffusion outputs: irregular borders, incorrect colors, and high chromatic noise.}
\label{fig:artifacts}
\end{figure}

This work proposes a digital processing pipeline to restore the structural integrity of these assets, reimposing the discretization that the probabilistic nature of diffusion neglects and reconciling the speed of synthetic generation with the precision required by digital asset engineering \cite{coutinho2024missing}. The study is guided by the premise that generative AI should not be interpreted as a direct replacement for the artist or for handcrafted pipelines, but as a rapid form-generation mechanism that requires a subsequent normalization stage before integration into technical environments. This aligns with the emerging paradigm of mixed-initiative co-creativity, where AI acts as a creative partner rather than an autonomous producer \cite{margarido2025boosting}. The article makes three main contributions. First, it formalizes the mismatch between diffusion-based generation and the discrete logic of Pixel Art as a domain-specific engineering problem rather than a purely visual mismatch. Second, it presents a deterministic post-processing architecture that restores grid alignment, reduces chromatic entropy, and normalizes transparency. Third, it discusses the distinction between technical compliance and stylistic authenticity, arguing that objective restoration is necessary but not sufficient to fully characterize genuine Pixel Art production. From these contributions, two research questions emerge naturally. \textbf{RQ1}: to what extent can deterministic image processing recover the technical properties that diffusion-based outputs systematically violate? \textbf{RQ2}: even when technical compliance is restored, how far can objective normalization approximate the perceptual expectations of expert Pixel Art practitioners? These questions motivate both the experimental design and the interpretive emphasis of the discussion.

\begin{figure}[ht]
\centering
\includegraphics[width=0.6\linewidth]{figures/latent_diffusion_process.png}
\caption{Latent diffusion process: resulting in interpolations and gradients incompatible with the discrete Pixel Art grid.}
\label{fig:difusion}
\end{figure}

\section{Theoretical Foundation}

The structural mismatch between Latent Diffusion Models (LDM) and Pixel Art is the theoretical axis of this work. While Pixel Art operates by discrete manipulation of pixels in low-resolution matrices \cite{chen2025gametilenet}, LDMs synthesize images through stochastic processes in continuous latent spaces, frequently generating interpolation artifacts and \textit{mixels} \cite{chen2023pixart}. Although approaches based on GANs \cite{coutinho2024pixel,kim2023game} and embeddings \cite{merino2023five} have advanced in sprite generation, the increasing adoption of diffusion models for character and animation sequences places this problem at the center of game asset development \cite{gallotta2025importance,ribeiro2024image,hsieh2025sprite}. The reverse sampling process introduces anti-aliasing that is technically prohibitive for the style: the validity of a Pixel Art asset requires hard edges and grid predictability \cite{mercioni2024innovative}, but AI generates thousands of tonal variations where an indexed and finite palette should exist. Similar challenges are observed in voxel art, where 3D meshes must be discretized into volumetric grids while preserving semantic coherence \cite{huang2026voxify3d}. To remedy this gap, this work resorts to Nearest-Neighbor Resampling, based on the preservation of high-frequency discontinuities, capable of reimposing the original geometry without introducing new gradients \cite{kopf2011,coutinho2024missing}. Chromatic reduction is treated via vector clustering, which synthesizes indexed palettes from the Euclidean color space \cite{coutinho2025imputation}. The theoretical framework establishes digital processing as a domain correction step for asset-oriented compliance criteria \cite{ekpete2023pixelperfect}.

\subsection{Pixel Art as a Constrained Representation}

Unlike general-purpose digital illustration, Pixel Art is built around intentional limitation. Resolution is low, pixel placement is semantically meaningful, and color selection is usually compact enough that each palette decision affects readability, shading, and silhouette. The visual language of the medium depends on preserving discrete transitions rather than smoothing them away. Therefore, any generative system intended for this domain must be evaluated not only by semantic coherence, but also by whether its output obeys low-level structural constraints such as palette compactness, edge sharpness, and grid consistency. This observation changes the usual success criterion for image generation. A diffusion model may correctly satisfy a textual prompt while still producing an unusable asset, because prompt alignment alone does not imply compliance with the discrete ontology of the medium. In other words, the central issue is not whether diffusion models can imitate the look of Pixel Art, but whether they can produce technically operational artifacts that behave like Pixel Art under real production constraints.

\subsection{Generation Quality Versus Asset Validity}

Recent literature on Pixel Art generation demonstrates growing interest in controllable asset creation through GANs, embeddings, and diffusion-based systems \cite{coutinho2022,coutinho2024pixel,kim2023game,merino2023five,saravanan2022}. Newer architectures such as PixelGen argue that end-to-end pixel diffusion can surpass latent-based methods by avoiding VAE bottlenecks, provided perceptual supervision is applied \cite{ma2026pixelgen}. Similarly, PixDiff-PIG introduces palette-informed diffusion frameworks that leverage OKLab color modeling to enforce chromatic consistency during the synthesis process \cite{li2025pixdiff}. Score distillation techniques like SD-$\pi$XL have also been proposed to produce low-resolution quantized imagery that maintains semantic faithfulness through differentiable generators \cite{binninger2024sdpixl}. However, most efforts prioritize image synthesis or stylistic transfer, not post-generation validation. This leaves a methodological gap between visually plausible outputs and deployable sprites. The gap is especially visible when generated images are resized, animated, or inserted into engines that assume crisp edges and indexed palettes. From this perspective, a valid Pixel Art workflow must distinguish between \emph{semantic generation quality} and \emph{asset validity}. Semantic quality refers to whether the generated image depicts the intended subject. Asset validity refers to whether the image satisfies the formal constraints required for use in production. The proposed pipeline addresses the second problem directly, complementing rather than competing with generative models.

\subsection{Why Automatic Validation Remains Difficult}

An important theoretical consequence of this work is that technical restoration does not fully solve the problem of authenticity. Some properties of Pixel Art can be measured objectively, such as palette size, edge discretization, or grid alignment. Others remain partially interpretive, including stylistic intentionality, readability under animation, and the degree to which a sprite reflects human compositional choices. This distinction is relevant because a processed image may become technically compliant while still lacking the visual economy or expressive judgment associated with handcrafted assets. Recent efforts to quantify mixed-initiative co-creativity attempt to bridge this gap by modeling human-AI collaboration dynamics, yet expert judgment remains the gold standard for final asset approval \cite{margarido2024miccy}. For that reason, this article treats validation as a layered concept. Objective metrics are appropriate for verifying structural restoration, but expert visual assessment remains important when discussing stylistic fidelity. The pipeline is thus framed as a normalization mechanism that makes generated content usable and inspectable, not as a final arbiter of artistic authenticity.

\subsection{Positioning Within the Literature}

The literature suggests at least three broad lines of research relevant to this article. The first investigates direct generation of Pixel Art-like content through adversarial or diffusion-based models \cite{coutinho2024pixel,kim2023game,chen2023pixart,serpa2019towards}. The second explores controllability, semantic guidance, or dataset design for game-related assets \cite{chen2025gametilenet,gallotta2025importance,merino2023five}. The third line focuses on style-preserving transformations and pixel-consistent rendering strategies, including methods that explicitly respect low-resolution visual structure or explore the innovation mechanisms of digital media art in the AI era \cite{coutinho2024missing,coutinho2025imputation,ekpete2023pixelperfect,ubaid2025realizing,liang2025research}.

The present work is positioned at the intersection of these lines, but with a different emphasis. Instead of asking how to generate Pixel Art from scratch, it asks how to make already-generated content technically valid. This distinction is important because production teams often adopt tools opportunistically: a generator may be useful even if imperfect, provided there is a reliable way to convert its output into an editable, engine-compatible intermediate representation. In this sense, the contribution of the article is less about replacing generation research and more about adding an engineering layer that makes such research operational in practice.

\begin{figure}[ht]
\centering
\includegraphics[width=0.6\linewidth]{figures/interpolation_vs_discretization.png}
\caption{Contrast between bilinear interpolation (left) and nearest-neighbor discretization (right)}
\label{fig:interp}
\end{figure}

\section{System Architecture and Implementation}

The proposed solution was built as a deterministic correction architecture composed of three sequential stages, designed to reverse the structural entropy inherent in diffusion-based outputs (Fig.~\ref{fig:pipeline_detailed}). This section details the design rationale, the mathematical formulation of each component, and the implementation decisions that ensure technical asset compliance.

\begin{figure}[ht]
\centering
\includegraphics[width=0.8\linewidth]{figures/restoration_pipeline_workflow.png}
\caption{Detailed restoration workflow: from raw misaligned AI output to the final normalized and engine-ready asset.}
\label{fig:pipeline_detailed}
\end{figure}

\subsection{Design Rationale and System Requirements}

The primary requirement for the pipeline is to restore the technical properties that diffusion models systematically violate: grid alignment, palette compactness, and transparency predictability. Instead of training a new generative model or relying on additional learned components, the system treats restoration as a deterministic normalization task. This decision has two major advantages: it ensures that the same pipeline remains compatible with different generators (DALL·E, MidJourney, Stable Diffusion) and it keeps the transformation stages fully inspectable and reproducible.

During the development of the broader project from which this paper is derived, a mapping of recurring failure modes in synthetic Pixel Art was conducted. The system was then designed to address these modes directly: softened edges are resolved through geometric reconstruction, excessive gradients are handled via chromatic quantization, and unstable alpha masks are normalized through binarization. This modular approach allows the pipeline to function as a technical validation layer between the stochastic generation phase and the engine-facing asset integration.

\subsection{Architectural Components and Mathematical Formulation}

The system architecture is organized as a sequential processing pipeline $P(I) = f_{\alpha}(f_{k}(f_{g}(I)))$, where $f_{g}$ is the geometric reconstruction, $f_{k}$ is the chromatic quantization, and $f_{\alpha}$ is the opacity binarization. This modularity ensures that each stage addresses a specific class of failure modes identified in diffusion outputs.

\textbf{1. Geometric Grid Reconstruction ($f_{g}$):}
To eliminate \textit{mixels} and restore the discrete grid topology, a Nearest-Neighbor (NN) mapping is applied. For an input image $I_{in} \in \mathbb{R}^{W_{in} \times H_{in} \times 4}$ and a target resolution $W_{out} \times H_{out}$, the mapping function $f_{g}: \mathbb{R}^2 \to \mathbb{R}^2$ is defined as:
\begin{equation}
I_{out}(x,y) = I_{in}\!\left(\left\lfloor x \cdot \frac{W_{in}}{W_{out}} + 0.5 \right\rfloor,\; \left\lfloor y \cdot \frac{H_{in}}{H_{out}} + 0.5 \right\rfloor\right)
\end{equation}
The implementation decision to use NN instead of Bilinear or Bicubic interpolation is critical: NN preserves the high-frequency edges and prevents the creation of intermediate color values, which is the primary source of 'blurry' artifacts in synthetic pixel art. Similar constant-color selection strategies have proven effective in broader image restoration contexts where channel integrity is paramount \cite{li2025scalable}. Furthermore, recent optimizations in diffusion training, including limited-interval guidance \cite{aila2024applying}, faster transformer architectures \cite{liu2024fasterdit,ge2024exploring}, and understanding objectives as evidence lower bounds \cite{gao2023understanding}, have accelerated the generation phase, making real-time post-processing even more critical for production responsiveness.

\textbf{2. Chromatic Quantization ($f_{k}$):}
The reduction of chromatic complexity is treated as a clustering problem in the RGB color space. The goal is to find a set of $K$ centroids $\{\mu_1, \dots, \mu_k\}$ that minimize the sum of squared distances from each pixel $p_i \in \mathbb{R}^3$ to its nearest centroid. While advanced vector quantization methods like Masked VQ offer robust latent compression \cite{nguyen2023masked}, our approach prioritizes a deterministic K-Means formulation as described in Algorithm~\ref{alg:kmeans} to ensure direct interpretability for artists.

\begin{algorithm}[ht]
\small
\caption{K-Means Color Quantization for Pixel Art}
\label{alg:kmeans}
\begin{algorithmic}[1]
\REQUIRE Image $I_{out} \in \mathbb{R}^{N \times 3}$, palette size $K$
\ENSURE Quantized Image $Q \in \mathbb{R}^{N \times 3}$, Palette $C$
\STATE Initialize $K$ centroids $\{\mu_1, \dots, \mu_k\}$ using $k$-means++
\REPEAT
    \STATE \textit{// Assignment step}
    \FOR{each pixel $p_i$ in $I_{out}$}
        \STATE $c^{(i)} \leftarrow \arg\min_{j} \| p_i - \mu_j \|^2$
    \ENDFOR
    \STATE \textit{// Update step}
    \FOR{each centroid $j \in \{1, \dots, k\}$}
        \STATE $\mu_j \leftarrow \frac{\sum_{i: c^{(i)}=j} p_i}{\sum_{i: c^{(i)}=j} 1}$
    \ENDFOR
\UNTIL{convergence (centroid displacement $< \epsilon$)}
\STATE Assign each pixel to its final centroid color
\end{algorithmic}
\end{algorithm}

\textbf{3. Opacity Binarization ($f_{\alpha}$):}
The final stage normalizes the alpha channel to satisfy engine-level binarization constraints. Given an input alpha channel $A_{in}$, the transformation is:
\begin{equation}
A_{out}(u,v) =
\begin{cases}
255 & \text{if } A_{in}(u,v) \geq \tau \\
0 & \text{if } A_{in}(u,v) < \tau
\end{cases}
\end{equation}
$\tau = 128$ was adopted based on technical sprite standards. This stage ensures that the resulting asset behaves predictably in engines that use 1-bit transparency masks for sorting and performance optimization.

\subsection{Computational Complexity and System Performance}

The overall complexity of the pipeline is dominated by the K-Means stage. For an image with $N$ pixels, $K$ colors, and $d=3$ dimensions (RGB), the complexity per iteration is $O(N \cdot K \cdot d)$. Given that $N$ is typically small for Pixel Art (e.g., $64 \times 64 = 4096$) and $K \leq 32$, the execution time is negligible (sub-millisecond) on modern hardware, allowing the pipeline to be used in real-time within interactive design tools. The geometric reconstruction and alpha binarization are $O(N)$, ensuring linear scalability with resolution.

\subsection{Prototype Implementation and Accessibility}

An important practical characteristic of the proposed pipeline is that it avoids dependence on retraining, prompt retuning, or generator-specific latent manipulations. The prototype was implemented using standard image-processing libraries, ensuring that all three stages are reproducible and can be integrated into broader asset management platforms. This is especially useful in research and production contexts where teams need stable, inspectable transformations rather than opaque learned corrections. The prototype demonstrates how classical image-processing methods can be operationalized in a lightweight environment that is easier to access and audit than complex end-to-end generative retraining workflows.

\begin{table}[ht]
\footnotesize
\centering
\begin{minipage}{0.48\linewidth}
\centering
\caption{Stage roles and trade-offs}
\label{tab:stages}
\resizebox{\linewidth}{!}{
\begin{tabular}{lp{2.8cm}p{2.5cm}}
\toprule
Stage & Role & Trade-off \\
\midrule
NN Resamp. & Reconstruct grid & Suppress softness \\
K-Means & Reduce entropy & Discard variation \\
Alpha Bin. & Norm. transparency & Remove semi-transp. \\
\bottomrule
\end{tabular}
}
\end{minipage}
\hfill
\begin{minipage}{0.48\linewidth}
\centering
\caption{Structured subset reported}
\label{tab:subset}
\resizebox{\linewidth}{!}{
\begin{tabular}{lp{4cm}}
\toprule
Aspect & Description \\
\midrule
Sources & DALL·E, MidJourney, SD \\
Subset & 120 curated images \\
Input Res. & Predom. 512$\times$512 \\
Output Res. & 64$\times$64 or 32$\times$32 \\
Inspection & 5 experts, 1--5 scale \\
\bottomrule
\end{tabular}
}
\end{minipage}
\end{table}

\subsection{Case Study and Evaluation Context}

The structured subset reported in the article inherits project-level constraints. It is a convenience sample rather than a benchmark dataset, and it was assembled to represent recurring artifact categories instead of balanced semantic classes. For that reason, the article treats the resulting measurements as descriptive evidence about the evaluated prototype behavior, not as definitive performance bounds for all generative image pipelines.

\subsection{Validation Protocol and Criteria}

Evaluation combined objective and interpretive criteria. Objective analysis focused on grid alignment, palette reduction, and alpha-channel normalization. In this article, ``technical compliance'' refers only to the conjunction of those selected structural criteria, not to a full certification of artistic quality or engine-readiness under every possible downstream use. Interpretive analysis was used to examine whether the processed images became more convincing as Pixel Art from the perspective of human specialists. The protocol reflects the broader research position of this article: technical correction and stylistic validation are related but not identical tasks. A processed asset can be mathematically cleaner than its raw counterpart and still benefit from human refinement. Rather than treating this as a weakness, the study interprets it as evidence that production-oriented Pixel Art pipelines should combine algorithmic normalization with expert inspection when authenticity is important.

\begin{table}[ht]
\footnotesize
\centering
\begin{minipage}{0.48\linewidth}
\centering
\caption{Evaluation dimensions}
\label{tab:criteria}
\resizebox{\linewidth}{!}{
\begin{tabular}{lp{4.5cm}}
\toprule
Dimension & Purpose \\
\midrule
Grid Alignment & Verify block convergence \\
Palette Comp. & Measure entropy reduction \\
Opacity Norm. & Confirm binary alpha \\
Visual Fidelity & Expert perception check \\
\bottomrule
\end{tabular}
}
\end{minipage}
\hfill
\begin{minipage}{0.48\linewidth}
\centering
\caption{Recurrent failure classes}
\label{tab:failures}
\resizebox{\linewidth}{!}{
\begin{tabular}{lp{5cm}}
\toprule
Failure class & Consequence for production \\
\midrule
Mixels & Unstable silhouettes \\
Gradients & Color proliferation \\
Soft Alpha & Mask ambiguity \\
Over-detail & Local visual noise \\
\bottomrule
\end{tabular}
}
\end{minipage}
\end{table}

\section{Software Design and System Architecture}

The proposed solution prioritizes modularity, inspectability, and ease of integration. Designed as a decoupled service, each normalization stage acts as an independent module. This architectural choice is inspired by the modular success of Scalable Interpolant Transformers (SiT) \cite{ma2024sit} and the efficiency of continuous latent generation models \cite{hang2025fast}. The prototype architecture is decomposed into three logical layers (Fig.~\ref{fig:arch_layers}): Ingestion, Normalization Core, and Export/Persistence. Implemented in Python, the system includes diagnostic logging to report modified pixel percentages and entropy reduction, ensuring transparency for technical artists. By incorporating structural enhancements like rotary position embeddings \cite{su2023roformer} in future ingestion modules, the system could better handle the spatial dependencies of uncropped sprite sheets. Its architectural skeleton supports future batch processing and temporal consistency modules, allowing it to function as a stylistic gatekeeper within studio pipelines.

\begin{figure}[ht]
\centering
% \includegraphics[width=0.6\linewidth]{figures/arch_layers.png}
\caption{Layered system architecture: ingestion, normalization, and export.}
\label{fig:arch_layers}
\end{figure}

\section{Experimental Evaluation and Results Analysis}

Evaluation was based on comparative quantitative analysis between raw assets and processed results. For the evaluated subset, inspection of block placement indicated that 87\% of original outputs exhibited grid deviations $> 0.5$ pixels. The mean positional error $E_p$ was eliminated by construction after geometric reconstruction (Fig.~\ref{fig:grid_error_details}). This addresses the 'reconstruction vs. generation' dilemma, where improving fidelity often conflicts with maintaining technical regularity in latent spaces \cite{yao2025reconstruction}.

\begin{figure}[ht]
\centering
\includegraphics[width=0.55\linewidth]{figures/grid_error_mixels_detail.png}
\caption{Visualizing the Grid Error ($E_p$): the transition from continuous color leakage (mixels) to a discrete grid-aligned structure.}
\label{fig:grid_error_details}
\end{figure}

Raw assets had an average of 35,400 colors; K-Means reduced chromatic complexity by $> 99.9\%$, while alpha analysis confirmed a pure bimodal distribution (0 or 255).

\subsection{Comparative Performance Analysis}

To establish the efficacy of the proposed architecture, we performed a comparative analysis against two common 'baseline' approaches used in image pixelization: (B1) standard Bicubic resampling followed by simple palette indexing, and (B2) Median-Cut quantization without prior grid reconstruction. The results (Table~\ref{tab:comparison}) indicate that while baseline B1 produces visually acceptable results for high-resolution images, it systematically fails to restore grid alignment in low-resolution assets. Baseline B2 tends to collapse subtle color transitions more aggressively than the K-Means stage, leading to loss of semantic detail in complex sprites.

\begin{table}[ht]
\footnotesize
\centering
\caption{Comparison against standard baselines}
\label{tab:comparison}
\resizebox{0.8\linewidth}{!}{
\begin{tabular}{lccc}
\toprule
Metric & Bicubic + Index (B1) & Median-Cut (B2) & Pipeline \\
\midrule
Grid Stability & Low & Medium & High \\
Chromatic Fidelity & Medium & Low & High \\
Alpha Mask Safety & Low & Low & High \\
Asset Validity & 12\% & 45\% & 100\% \\
\bottomrule
\end{tabular}
}
\end{table}

The proposed pipeline's superiority stems from the \textit{ordered sequentiality} of its modules. By reconstructing the grid before chromatic quantization, the system ensures that the K-Means algorithm operates on spatially stabilized blocks, preventing the clustering of anti-aliased gradients that would otherwise contaminate the final palette.

\subsection{Ablation Study: The Impact of Modular Sequentiality}

An ablation study confirmed that technical asset validity is an emergent property of the integrated pipeline. Removing Stage 1 (NN grid reconstruction) led to a 15\% increase in residual chromatic entropy, as clustering attempted to model interpolation noise. Bypassing Stage 3 (Alpha Binarization) resulted in silhouette artifacts during engine import. These findings validate the necessity of the complete architectural chain.

\subsection{Synthesis and Engineering Implications}

The evidence collected leads to three major engineering implications. First, the system reduces technical debt by automating grid alignment and palette indexing, preventing structural inconsistencies. Second, it enables scalable workflows where bulk work is performed in constant time by the normalization core. Third, architectural decoupling allows it to function as a stylistic gatekeeper, preserving visual coherence. Results indicate that AI generation time savings are better leveraged when integrated with a network-ready normalization layer.

\subsection{Limitations, threats, and research agenda}

Despite restoration, limitations remain: palette reduction can suppress detail in dense textures, and the method does not yet guarantee consistency across animation frames. While improving technical compliance, artistic authenticity still depends on human interpretation. Threats include the convenience nature of the corpus and the limited number of specialists. Even so, deterministic normalization substantially reduced the structural gap in diffusion outputs, opening a bridge for co-designing generation and normalization modules.

The research agenda points toward set-aware normalization for character poses and the integration of assessment modules to estimate perceptual conviction, ideally integrated directly into authoring tools. Future developments in narrative-driven scene generation \cite{chen2025narrative}, unified image tokenization \cite{qu2024tokenflow,zha2024language}, and cluster-oriented token prediction \cite{hu2025improving} may provide the necessary semantic grounding to automate these refinements. Furthermore, exploring decoupled diffusion transformers \cite{wang2025ddt}, pixel neural field diffusion \cite{wang2025pixnerd}, and native-resolution synthesis \cite{wang2025native} could minimize the initial entropy of AI outputs, reducing the corrective burden on the post-processing layer.


\section{Selected Case Studies and Visual Evidence}

To demonstrate practical utility, we analyze a character sprite (Medieval Knight) and a forest landscape background. The knight, generated via Stable Diffusion, exhibited severe chromatic noise; normalization produced a 64$\times$64 asset with 16 colors, eliminating over 30,000 redundant values and making it engine-ready. The forest scene, from MidJourney, contained over 50,000 colors; the pipeline collapsed this into 32 colors while preserving hue clusters and retro aesthetics. These cases show that the system recovers the 'economic' nature of Pixel Art, where pixel placement is intentional. While robust, the restoration involves trade-offs: extreme detail may be merged if Euclidean-close in RGB space, reinforcing the pipeline's role as an assistive layer where artists focus effort on final stylistic refinement.

\begin{figure}[ht]
\centering
\begin{minipage}{0.48\linewidth}
  \centering
  \includegraphics[width=\linewidth]{figures/case_study_landscape_raw.png}
\end{minipage}
\hfill
\begin{minipage}{0.48\linewidth}
  \centering
  \includegraphics[width=\linewidth]{figures/case_study_landscape_fixed.png}
\end{minipage}
\caption{Comparison: Raw diffusion output (left) vs processed asset (right)}
\label{fig:comparativo}
\end{figure}

\begin{table}[ht]
\footnotesize
\centering
\begin{minipage}{0.48\linewidth}
\centering
\caption{Restoration Metrics}
\label{tab:res}
\resizebox{\linewidth}{!}{
\begin{tabular}{lcc}
\toprule
Metric & Raw & Pipeline \\
\midrule
Colors & 35,400 & 16-32 \\
Grid Err & 0.54 & 0.00 \\
Alpha Var & 124.2 & 0.00 \\
Compl. & 0\% & 100\% \\
\bottomrule
\end{tabular}
}
\end{minipage}
\hfill
\begin{minipage}{0.48\linewidth}
\centering
\caption{Specialist eval.}
\label{tab:subjective}
\resizebox{\linewidth}{!}{
\begin{tabular}{lcc}
\toprule
Criterion & Raw & Proc. \\
\midrule
Fidelity & 2.3 & 4.1 \\
Coherence & Low & High \\
Readability & Low & High \\
Usability & Lim. & Cons. \\
\bottomrule
\end{tabular}
}
\end{minipage}
\end{table}

\section{Conclusion}

The investigation conducted in this work indicates that latent diffusion models do not automatically translate into technical viability for Pixel Art. Quantitative evidence shows that raw outputs often remain structurally unsuitable for asset-oriented workflows. The proposed pipeline proved effective: geometric reconstruction, K-Means quantization, and Alpha binarization reduced chromatic complexity by more than 99.9\%, transforming aesthetic references into functional assets. Technical restoration and artistic authenticity should not be conflated; the pipeline reliably enforces structural compliance and improves perceived fidelity, but final legitimacy depends on intentional design choices. We conclude that the value of generative AI in this niche lies in its utility as a form generation engine, conditioned on a rigorous technical normalization layer. Future work should advance toward consistency-aware processing for animation frames and hybrid validation schemes combining objective metrics with expert-informed models of stylistic authenticity.

\begin{credits}
\subsubsection{\ackname} The authors thank the UFMT for the institutional support provided during this research.
\end{credits}

%
% ---- Bibliography ----
%
\footnotesize
\bibliographystyle{splncs04}
\bibliography{article-bracis}

\end{document}
